\section{The full correlations of the original relation rows}
\label{sec:pgrc}

Retain the complete quartet, full Taylor unit, original theta measure,
ordered tensor coordinates and monic relation columns of PGRT.1--8.
In particular $k,m\geq1$, $0<\delta<1/2$, $\gamma>2$,
$q=[1+k(m-1)](k+1)^2$, and $E=\mathbb C[S]/(\chi)$ has its
original remainder frame $1,S,\ldots,S^{q-1}$. Every $a_j>0$ is
the original squared length of $d_j=\chi p_{j+1-q}$, and
$\alpha_j:E\to\mathbb C$ is the exact integral PGRT.8. No relation
column or measure is rescaled below. For finite coefficient matrices,
stars denote conjugate transposes in the indicated frames. For maps
into a stated Hilbert target, including $Z$ below, stars denote
adjoints for that target and the stated coefficient metric. All
Hilbert inner products are conjugate linear in the first variable.

\subsection{A finite source map for all the retained correlations}

Fix the original last endpoint and the complete row index set
\[
 H=\widehat G_{2q},\qquad I=\{q-1,q,\ldots,2q-1\},\qquad
 C_{ij}=\alpha_iH^{-1}\alpha_j^*\quad(i,j\in I).
 \tag{PGRC.1}\label{eq:pgrc-original-gram}
\]
Here $H>0$ is the actual graph quotient form computed by PGS.9--13.
Thus each entry of $C$ is determined by the original finite theta
moments, the full tensor remainder, and the inverse of that form.
For $J\subset I$, in the inherited order, let $A_J$ have rows
$\alpha_j$, let $D_J=\operatorname{diag}(a_j:j\in J)$ and put
\[
 H_J=H+A_J^*D_J^{-1}A_J,
 \qquad B_J=D_J+C_{JJ}.
 \tag{PGRC.2}\label{eq:pgrc-subset-forms}
\]
The empty set means the zero-dimensional coordinate space, zero
row map, and $H_\varnothing=H$. The determinants of
$B_\varnothing$ and $D_\varnothing$ are one, while
$\det H_\varnothing=\det H$. Both
$H_J$ and $B_J$ are positive definite on their respective spaces:
their quadratic forms are the sum of a positive original form
and a nonnegative squared norm.

The precise graph realization, including its inverse on the image,
is
\[
 \begin{aligned}
 V_J:(E,H_J)&\longrightarrow(E,H)\oplus(\mathbb C^J,D_J^{-1}),
 &V_Jx&=(x,A_Jx),\\
 V_J^\dagger(y,z)&=H_J^{-1}(Hy+A_J^*D_J^{-1}z),
 &V_J^\dagger V_J&=I_E,\\
 \ker V_J&=\{0\},
 &\ker V_J^\dagger&=\{(y,z):Hy+A_J^*D_J^{-1}z=0\}.
 \end{aligned}
 \tag{PGRC.3}\label{eq:pgrc-source-map}
\]
Indeed expanding the norm of $(x,A_Jx)$ gives $x^*H_Jx$,
and pairing this vector with $(y,z)$ proves the stated adjoint.
Multiplication proves the left inverse and the kernel formula.
Consequently $V_JV_J^\dagger$ is the orthogonal projection onto
this image; it is self-adjoint and its square equals itself.
This construction retains each original relation length in the
coordinate-space metric $D_J^{-1}$.

Direct multiplication, with every factor in its stated order, gives
\[
 H_J^{-1}=H^{-1}-H^{-1}A_J^*B_J^{-1}A_JH^{-1}.
 \tag{PGRC.4}\label{eq:pgrc-subset-inverse}
\]
To verify it, multiply the right side on the left by $H_J$.
The extra terms, after factoring $A_J^*$ on the left and
$A_JH^{-1}$ on the right, have middle factor
$D_J^{-1}-B_J^{-1}-D_J^{-1}C_{JJ}B_J^{-1}=0$,
because $I+D_J^{-1}C_{JJ}=D_J^{-1}B_J$.
The positive forms are invertible, so this is the claimed inverse.

For $j\in I$ and a subset of its strictly later rows
$J\subset J_j:=\{j+1,\ldots,2q-1\}$ define
\[
 r_j(J)=C_{jJ}B_J^{-1}C_{Jj},\qquad
 v_j(J)=\frac{C_{jj}-r_j(J)}{a_j}
       =\frac{\alpha_jH_J^{-1}\alpha_j^*}{a_j}\geq0.
 \tag{PGRC.5}\label{eq:pgrc-correlation-removal}
\]
The empty-set correction is zero. The second expression follows
by PGRC.4 and proves the sign without discarding a cross term.
The correction $r_j(J)$ is zero exactly when $C_{Jj}=0$, since
$B_J^{-1}>0$. In particular its strictness is determined by the
actual original-metric correlations of the rows, not by their
mere presence as formal coordinates.

PGRT.6 gives the precise identification for all later rows:
\[
 H_{J_j}=\widehat G_{j+1},\qquad
 \eta_j=v_j(J_j)
  =\frac{C_{jj}-C_{jJ_j}(D_{J_j}+C_{J_jJ_j})^{-1}C_{J_jj}}{a_j}.
 \tag{PGRC.6}\label{eq:pgrc-exact-row}
\]
The first identity telescopes the finite sum
$\widehat G_{j+1}-\widehat G_{2q}
=\sum_{i=j+1}^{2q-1}\alpha_i^*\alpha_i/a_i$.
Inserting it into the proved inverse calculation gives the second.
This evaluates exactly the same coefficient $\eta_j$ as PGRT.12.

\subsection{Every omitted correlation and an explicit finite remainder}

For disjoint $J,K\subset J_j$, define the full remaining block
and row, retaining their subtraction:
\[
 \begin{aligned}
 S_{K\mid J}&=D_K+C_{KK}-C_{KJ}B_J^{-1}C_{JK},\\
 e_{j,K\mid J}&=C_{jK}-C_{jJ}B_J^{-1}C_{JK}.
 \end{aligned}
 \tag{PGRC.7}\label{eq:pgrc-remaining-block}
\]
They have exact source expressions
\[
 S_{K\mid J}=D_K+A_KH_J^{-1}A_K^*\succeq D_K>0,
 \qquad e_{j,K\mid J}=\alpha_jH_J^{-1}A_K^*.
 \tag{PGRC.8}\label{eq:pgrc-remaining-source}
\]
These follow by substituting PGRC.4. The block calculation uses
the exact coordinate restriction permutation
$P_{J,K}:\mathbb C^{J\cup K}\to\mathbb C^J\oplus\mathbb C^K$,
$z\mapsto(z|_J,z|_K)$, whose inverse reassembles the coordinates
in the original inherited order. Thus
$P_{J,K}B_{J\cup K}P_{J,K}^*$ has diagonal blocks
$B_J,D_K+C_{KK}$ and off-diagonal blocks $C_{JK},C_{KJ}$.
The row $C_{j,J\cup K}P_{J,K}^*$ is $(C_{jJ},C_{jK})$.
Since $P_{J,K}^*P_{J,K}=I$, their inverse quadratic pairing
$r_j(J\cup K)$ is unchanged. This explicitly covers interleaved
indices. Block elimination in these specified coordinates gives
\[
 r_j(J\cup K)-r_j(J)
 =e_{j,K\mid J}S_{K\mid J}^{-1}e_{j,K\mid J}^*,
 \tag{PGRC.9}\label{eq:pgrc-exact-remainder}
\]
which also holds if either block is empty. Here is the elimination
calculation. The equations
$B_Jx+C_{JK}y=C_{Jj}$ and
$C_{KJ}x+(D_K+C_{KK})y=C_{Kj}$ first give
$x=B_J^{-1}(C_{Jj}-C_{JK}y)$ and then
$S_{K\mid J}y=e_{j,K\mid J}^*$. Pair the first solution and this
second solution with $(C_{jJ},C_{jK})$; the resulting scalar is
$r_j(J)+e_{j,K\mid J}S_{K\mid J}^{-1}e_{j,K\mid J}^*$,
which is the left side's full block inverse expression.

Since $S_{K\mid J}\succeq D_K$, inverse order yields the fully
explicit estimate
\[
 0\leq r_j(J\cup K)-r_j(J)
 \leq\sum_{i\in K}\frac{|(e_{j,K\mid J})_i|^2}{a_i}.
 \tag{PGRC.10}\label{eq:pgrc-remainder-bound}
\]
For completeness, conjugating the positive inequality by
$D_K^{-1/2}$ produces a positive matrix with eigenvalues at least
one. Its inverse has eigenvalues at most one; undoing the
congruence proves $S_{K\mid J}^{-1}\preceq D_K^{-1}$.
Pairing with the original row proves PGRC.10. No row's length
has changed during this argument.

For a single additional index $b$, PGRC.9 is the exact update
\[
 r_j(J\cup\{b\})=r_j(J)+
 \frac{|C_{jb}-C_{jJ}B_J^{-1}C_{Jb}|^2}
 {a_b+C_{bb}-C_{bJ}B_J^{-1}C_{Jb}}.
 \tag{PGRC.11}\label{eq:pgrc-single-update}
\]
Its denominator is at least $a_b>0$ by PGRC.8. The update
is strictly positive exactly when its displayed numerator is
nonzero. Thus ordered row insertion monotonically decreases
$v_j(J)$; every equality case is given by the stated cross row.

Choose any subsets $J(j)\subset J_j$, and let
$K(j)=J_j\setminus J(j)$. Set
\[
 \begin{aligned}
 \overline\eta_j&=v_j(J(j)),\\
 \underline\eta_j
 &=\max\left\{l_j,\ v_j(J(j))-
       \sum_{i\in K(j)}
       \frac{|(e_{j,K(j)\mid J(j)})_i|^2}{a_ja_i}\right\},
 \quad l_j=\frac{\alpha_j\widehat G_{q-1}^{-1}\alpha_j^*}{a_j}.
 \end{aligned}
 \tag{PGRC.12}\label{eq:pgrc-row-enclosure}
\]
PGRT.12 gives $0\leq l_j\leq\eta_j$. Equations PGRC.9--10
with $J\cup K=J_j$ prove
$0\leq\underline\eta_j\leq\eta_j\leq\overline\eta_j$.
Using the original weights
$\omega_{q-1}=\omega_{2q-1}=1$ and $\omega_j=2$ in between,
one obtains the finite original-source bounds
\[
 \begin{aligned}
 L_k^{\rm corr}(J)&=\sum_{j=q-1}^{2q-1}
                         \omega_j\log(1+\underline\eta_j),\\
 U_k^{\rm corr}(J)&=\sum_{j=q-1}^{2q-1}
                         \omega_j\log(1+\overline\eta_j),\\
 L_k^{\rm rel}&\leq L_k^{\rm corr}(J)
 \leq\widehat{\mathcal B}_k\leq U_k^{\rm corr}(J)
 \leq U_{k,2q}^{\rm rel}.
 \end{aligned}
 \tag{PGRC.13}\label{eq:pgrc-volume-enclosure}
\]
The first inequality uses $\underline\eta_j\geq l_j$;
the last uses $r_j(J)\geq0$ and
$C_{jj}/a_j=u_{j,2q}$ in PGRT.11. The middle inequalities
follow from the increasing logarithm and the exact weighted sum
PGRT.12. Taking all $J(j)=J_j$ makes both bounds equal to
$\widehat{\mathcal B}_k$, since then every remainder is zero.
Taking the first $\min(n,|J_j|)$ later rows gives decreasing
upper bounds, exact after at most $q$ insertions per row. Each
step has the original finite dimensions and can be evaluated
using PGRC.11. Lower bounds may be retained by taking the
maximum over all previous insertions, which preserves their
validity and makes that retained sequence increasing.

The excess of the partial upper bound is itself evaluated by
\[
 U_k^{\rm corr}(J)-\widehat{\mathcal B}_k
 =\sum_j\omega_j\log\left(1+
 \frac{e_{j,K(j)\mid J(j)}S_{K(j)\mid J(j)}^{-1}
                 e_{j,K(j)\mid J(j)}^*}
 {a_j(1+\eta_j)}\right).
 \tag{PGRC.14}\label{eq:pgrc-volume-error}
\]
Indeed PGRC.9 gives
$\overline\eta_j-\eta_j=eS^{-1}e^*/a_j$; divide
$(1+\overline\eta_j)/(1+\eta_j)$ and take logarithms.
This also proves a finite computable bound by replacing the
numerator with PGRC.10 and $\eta_j$ with
$\underline\eta_j$, retaining each sign and each row length.

\subsection{Both negative corrections in the same original inverse}

Keep $Z=Z_{2q}^{\rm mix}=(D_{\rm mix},f_{2q})$ from PGRT.9,
mapping into its original ordered Hilbert sum. PGRT.10--11 and
PGRC.6 together give the exact double subtraction
\[
 \begin{aligned}
 a_j\eta_j={}&\alpha_jK_{\rm low}^{-1}\alpha_j^*\\
 &-\left\|(I+ZK_{\rm low}^{-1}Z^*)^{-1/2}
                      ZK_{\rm low}^{-1}\alpha_j^*\right\|^2\\
 &-C_{jJ_j}(D_{J_j}+C_{J_jJ_j})^{-1}C_{J_jj}.
 \end{aligned}
 \tag{PGRC.15}\label{eq:pgrc-double-subtraction}
\]
Every $C$ on the last line is computed using the same full
$H=\widehat G_{2q}$; it is not replaced by a Gram in
$K_{\rm low}$. The first subtracted squared norm is positive
for $k\geq2$ and $\alpha_j\ne0$, by the proved injectivity of
$D_{\rm mix}$. The second is positive exactly when
$C_{J_jj}\ne0$, by PGRC.5. Both can vanish in their explicitly
stated cases. Their sum remains at most the first positive term,
because the left side is the positive inverse quadratic form
of PGRC.6. Thus the mixed residual and the relation correlations
have been carried into one original quantity with their full
negative signs.

The arithmetic receiving calculation retains the exact graph
route HC9g. For integers $p,s\geq0$ use the proved ACM27
$g_s^\pm=g_{s,\rm gr}^\pm$ and define
\[
 \begin{aligned}
 b_{p,s}^{\rm corr,-}
 &=\max\{\beta_{p,\Gamma}^-,L_k^{\rm corr}(J)-g_s^+\},\\
 b_{p,s}^{\rm corr,+}
 &=\min\{\beta_{p,\Gamma}^+,U_k^{\rm corr}(J)-g_s^-\}.
 \end{aligned}
 \tag{PGRC.16}\label{eq:pgrc-arithmetic-return}
\]
Subtraction of the two proved bounds on the literal signed
augmentation gives an interval containing $\mathcal B_k$.
The endpoints obey the exact inclusion order
\[
 [b_{p,s}^{\rm gr,-},b_{p,s}^{\rm gr,+}]
 \subset [b_{p,s}^{\rm corr,-},b_{p,s}^{\rm corr,+}]
 \subset [b_{p,s}^{\rm rel,-},b_{p,s}^{\rm rel,+}]
 \subset [\beta_{p,\Gamma}^-,\beta_{p,\Gamma}^+].
 \tag{PGRC.17}\label{eq:pgrc-nesting}
\]
For the first inclusion use
$L_k^{\rm corr}\leq\widehat{\mathcal B}_k\leq U_k^{\rm corr}$;
for the second use PGRC.13; and for the third use the maximum
and minimum in the definitions. With all later rows included,
the first two endpoints equal the exact graph endpoints HC9g.
The computed interval width is at most
\[
 \min\{2e_{p,\Gamma},\,
       U_k^{\rm corr}(J)-L_k^{\rm corr}(J)+2e_{s,\rm gr}\},
 \tag{PGRC.18}\label{eq:pgrc-arithmetic-width}
\]
because ACM16 and ACM27 place their endpoints inside the stated
centered intervals. Intersection cannot increase either width.

For HC9--10 and HC14g preserve their domain $k\geq3$, the original
$0<b<\pi/2$, $\gamma>2$, and the phase and circumference comparison
specified in PGRT.14's receiving paragraph. Subtract exactly
$T_k=4q\log(D_hk)$ and $\mathcal B_k^\Gamma$ for the respective
residuals, intersect with their retained earlier endpoints, and
insert the arithmetic upper value into the three increasing
phase functions. This proves the same exact receiving maps as
PGRT.13--14 with the more precise finite row enclosure PGRC.16.
It leaves the full exact-graph path available, rather than
replacing it by a weaker evaluation.

All maps in PGRC.3 and all inverse-image kernel descriptions act
coordinatewise on a fixed original coefficient mask. The empty
present mask retains its outer label, external absence retains
its distinct label and singleton inverse image, and every
represented-zero fibre retains the displayed linear kernel in
each active coordinate. The full finite calculation has at
most $q+1$ relation rows with $q=[1+k(m-1)](k+1)^2$; these
dimensions remain in the source integrals and inverse blocks.
No fixed-packet row convergence has been used to assert an
unevaluated uniform bound in the separate variable $k$.
